% ============================================================
\chapter{Le Mod\`ele Entit\'e-Relation}
\label{ch:entite_relation}
% ============================================================

Avant d'\'ecrire la premi\`ere ligne de SQL, il faut \emph{penser} la
structure des donn\'ees. C'est l'\'etape de la \emph{mod\'elisation
conceptuelle}, et l'outil le plus utilis\'e pour cela est le mod\`ele
entit\'e-relation, propos\'e par Peter Chen en 1976. L'id\'ee est de
repr\'esenter le monde r\'eel par des \emph{entit\'es} (les objets), des
\emph{attributs} (leurs propri\'et\'es) et des \emph{relations} (les liens
entre eux), dans un diagramme visuel qui sert de plan d'architecte pour la
base de donn\'ees.

% ------------------------------------------------------------
\section{Introduction \`a la conception de bases de donn\'ees}
% ------------------------------------------------------------

La conception d'une base de donn\'ees suit trois \'etapes principales~:
\begin{enumerate}
    \item \textbf{Analyse des besoins} --- recueil des exigences aupr\`es des utilisateurs.
    \item \textbf{Conception conceptuelle} --- mod\'elisation ind\'ependante de tout SGBD
          (mod\`ele Entit\'e-Relation).
    \item \textbf{Conception logique} --- traduction en sch\'ema relationnel (tables SQL).
    \item \textbf{Conception physique} --- choix des index, du partitionnement, etc.
\end{enumerate}

Le mod\`ele Entit\'e-Relation (E/R), introduit par Peter Chen en 1976, est
l'outil standard pour la conception conceptuelle.

\section{Concepts fondamentaux}

\begin{definition}[Entit\'e]
Une \emph{entit\'e} est un objet du monde r\'eel, distinguable des autres objets,
sur lequel on souhaite stocker des informations. Un \emph{type d'entit\'e}
regroupe les entit\'es partageant les m\^emes attributs.
\end{definition}

\begin{definition}[Attribut]
Un \emph{attribut} est une propri\'et\'e qui d\'ecrit une entit\'e ou une relation.
Chaque attribut a un \emph{domaine} (ensemble des valeurs possibles).
On distingue~:
\begin{itemize}
    \item \textbf{Attribut simple} vs.\ \textbf{compos\'e} (ex.\ adresse = rue + ville + code postal).
    \item \textbf{Attribut mono-valu\'e} vs.\ \textbf{multi-valu\'e} (ex.\ t\'el\'ephones).
    \item \textbf{Attribut stock\'e} vs.\ \textbf{d\'eriv\'e} (ex.\ \^age calcul\'e \`a partir de la date de naissance).
\end{itemize}
\end{definition}

\begin{definition}[Cl\'e d'entit\'e]
Un attribut (ou ensemble d'attributs) dont la valeur identifie de mani\`ere
unique chaque instance d'un type d'entit\'e. On le souligne dans le diagramme E/R.
\end{definition}

\begin{definition}[Relation (association)]
Une \emph{relation} (ou \emph{association}) repr\'esente un lien entre deux ou
plusieurs types d'entit\'es. Elle peut poss\'eder des attributs propres.
\end{definition}

\section{Cardinalit\'es}

\begin{definition}[Cardinalit\'e]
La \emph{cardinalit\'e} d'une participation d'un type d'entit\'e \`a une relation
exprime le nombre minimal et maximal d'instances de la relation auxquelles
une instance de l'entit\'e peut participer. On note $(min, max)$.
\end{definition}

Les cardinalit\'es les plus courantes sont~:

\begin{formules}[title=\textbf{Types de cardinalit\'es}]
\begin{center}
\begin{tabular}{lll}
\toprule
\textbf{Notation} & \textbf{Signification} & \textbf{Exemple} \\
\midrule
$1:1$ (un-\`a-un) & Chaque entit\'e est li\'ee \`a au plus une & Personne --- Passeport \\
$1:N$ (un-\`a-plusieurs) & Une entit\'e est li\'ee \`a plusieurs & Professeur --- Cours \\
$N:M$ (plusieurs-\`a-plusieurs) & Plusieurs \`a plusieurs & \'Etudiant --- Cours \\
\bottomrule
\end{tabular}
\end{center}
\end{formules}

\begin{exemple}
Un professeur enseigne \emph{plusieurs} cours, mais chaque cours est enseign\'e
par \emph{un seul} professeur~: c'est une relation $1:N$.

Un \'etudiant peut s'inscrire \`a \emph{plusieurs} cours, et chaque cours accueille
\emph{plusieurs} \'etudiants~: c'est une relation $N:M$.
\end{exemple}

\section{Diagrammes E/R avec TikZ}

\begin{center}
\begin{tikzpicture}[
    entity/.style={rectangle, draw, fill=blue!10, minimum width=2.5cm, minimum height=1cm, font=\bfseries},
    relationship/.style={diamond, draw, fill=orange!20, minimum width=2cm, minimum height=1cm, aspect=2, font=\small\bfseries},
    attribute/.style={ellipse, draw, fill=green!10, font=\small},
    keyattr/.style={ellipse, draw, fill=yellow!20, font=\small\bfseries},
    every edge/.style={draw, -}
]
    % Entites
    \node[entity] (etudiant) at (0, 0) {\'Etudiant};
    \node[entity] (cours)    at (8, 0) {Cours};
    \node[entity] (prof)     at (8, -4) {Professeur};

    % Relations
    \node[relationship] (inscrit) at (4, 0) {Inscrit};
    \node[relationship] (enseigne) at (8, -2) {Enseigne};

    % Liens entite-relation
    \draw (etudiant) -- node[above] {$N$} (inscrit);
    \draw (inscrit) -- node[above] {$M$} (cours);
    \draw (prof) -- node[right] {$1$} (enseigne);
    \draw (enseigne) -- node[right] {$N$} (cours);

    % Attributs Etudiant
    \node[keyattr] (eid) at (-2, 1.5) {id\_etud};
    \node[attribute] (enom) at (0, 1.8) {nom};
    \node[attribute] (epre) at (2, 1.5) {pr\'enom};

    \draw (etudiant) -- (eid);
    \draw (etudiant) -- (enom);
    \draw (etudiant) -- (epre);

    % Attributs Cours
    \node[keyattr] (cid) at (6, 1.5) {id\_cours};
    \node[attribute] (ctit) at (8, 1.8) {intitul\'e};
    \node[attribute] (ccred) at (10, 1.5) {cr\'edits};

    \draw (cours) -- (cid);
    \draw (cours) -- (ctit);
    \draw (cours) -- (ccred);

    % Attribut de la relation Inscrit
    \node[attribute] (note) at (4, 1.5) {note};
    \draw (inscrit) -- (note);

    % Attributs Professeur
    \node[keyattr] (pid) at (6, -5) {id\_prof};
    \node[attribute] (pnom) at (8, -5.5) {nom};
    \node[attribute] (pdep) at (10, -5) {d\'ept};

    \draw (prof) -- (pid);
    \draw (prof) -- (pnom);
    \draw (prof) -- (pdep);
\end{tikzpicture}
\end{center}

\section{Entit\'es faibles}

\begin{definition}[Entit\'e faible]
Une \emph{entit\'e faible} est un type d'entit\'e qui ne poss\`ede pas de cl\'e propre
suffisante pour l'identifier de mani\`ere unique. Son identification d\'epend
d'un type d'entit\'e \emph{propri\'etaire} (entit\'e forte) via une
\emph{relation identifiante}.
\end{definition}

\begin{exemple}
Consid\'erons les salles de cours d'un b\^atiment. Une salle est identifi\'ee
par son num\'ero \emph{dans} un b\^atiment donn\'e. Le type d'entit\'e \texttt{Salle}
est faible, identifi\'e par la combinaison de \texttt{Batiment.id\_bat} et
\texttt{Salle.numero}.
\end{exemple}

\begin{center}
\begin{tikzpicture}[
    entity/.style={rectangle, draw, fill=blue!10, minimum width=2.5cm, minimum height=1cm, font=\bfseries},
    weakentity/.style={rectangle, draw, double, fill=blue!10, minimum width=2.5cm, minimum height=1cm, font=\bfseries},
    relationship/.style={diamond, draw, fill=orange!20, minimum width=1.8cm, minimum height=0.9cm, aspect=2, font=\small\bfseries},
    weakrel/.style={diamond, draw, double, fill=orange!20, minimum width=1.8cm, minimum height=0.9cm, aspect=2, font=\small\bfseries},
]
    \node[entity] (bat) at (0,0) {B\^atiment};
    \node[weakrel] (contient) at (4,0) {Contient};
    \node[weakentity] (salle) at (8,0) {Salle};

    \draw (bat) -- node[above] {$1$} (contient);
    \draw[double] (contient) -- node[above] {$N$} (salle);
\end{tikzpicture}
\end{center}

\section{Sp\'ecialisation et g\'en\'eralisation}

\begin{definition}[Sp\'ecialisation / G\'en\'eralisation]
La \emph{sp\'ecialisation} consiste \`a d\'efinir des sous-types d'un type d'entit\'e
g\'en\'eral (relation \og est-un \fg, \emph{is-a}). La \emph{g\'en\'eralisation}
est le processus inverse~: regrouper des types d'entit\'es partageant des
attributs communs en un super-type.
\end{definition}

\begin{exemple}
Le type \texttt{Personne} peut \^etre sp\'ecialis\'e en \texttt{\'Etudiant} et
\texttt{Professeur}. Les attributs communs (nom, pr\'enom) sont dans
\texttt{Personne}, les attributs sp\'ecifiques (fili\`ere pour \'Etudiant,
d\'epartement pour Professeur) sont dans les sous-types.
\end{exemple}

Contraintes de sp\'ecialisation~:
\begin{itemize}
    \item \textbf{Disjointe} (d) vs.\ \textbf{chevauchante} (o)~: une entit\'e
          peut-elle appartenir \`a plusieurs sous-types ?
    \item \textbf{Totale} (t) vs.\ \textbf{partielle} (p)~: toute entit\'e
          du super-type doit-elle appartenir \`a au moins un sous-type ?
\end{itemize}

\section{Passage du mod\`ele E/R au mod\`ele relationnel}

\begin{theoreme}[R\`egles de traduction E/R $\to$ Relationnel]
Les r\`egles de passage sont les suivantes~:
\begin{enumerate}
    \item \textbf{Type d'entit\'e fort} $\to$ une table dont la cl\'e primaire est
          la cl\'e de l'entit\'e.
    \item \textbf{Type d'entit\'e faible} $\to$ une table dont la cl\'e primaire
          combine la cl\'e partielle et la cl\'e de l'entit\'e propri\'etaire
          (cl\'e \'etrang\`ere).
    \item \textbf{Relation $1:N$} $\to$ ajout d'une cl\'e \'etrang\`ere dans la table
          c\^ot\'e $N$.
    \item \textbf{Relation $N:M$} $\to$ cr\'eation d'une table d'association dont la
          cl\'e primaire est la paire des cl\'es \'etrang\`eres.
    \item \textbf{Relation $1:1$} $\to$ cl\'e \'etrang\`ere dans l'une des deux tables
          (de pr\'ef\'erence celle avec participation totale).
    \item \textbf{Attribut multi-valu\'e} $\to$ table s\'epar\'ee avec cl\'e \'etrang\`ere.
    \item \textbf{Sp\'ecialisation} $\to$ plusieurs strat\'egies possibles
          (tables s\'epar\'ees, table unique avec discriminateur, etc.).
\end{enumerate}
\end{theoreme}

\begin{exemple}
Appliquons les r\`egles au diagramme pr\'ec\'edent~:

\begin{codeblock}[title=\textbf{Traduction du diagramme E/R en SQL}]
\begin{minted}{sql}
-- Regle 1 : entite forte -> table
CREATE TABLE Etudiants (
    id_etudiant INTEGER PRIMARY KEY,
    nom         TEXT NOT NULL,
    prenom      TEXT NOT NULL
);

CREATE TABLE Cours (
    id_cours    INTEGER PRIMARY KEY,
    intitule    TEXT NOT NULL,
    credits     INTEGER DEFAULT 3,
    id_prof     INTEGER REFERENCES Professeurs(id_prof)  -- Regle 3 : 1:N
);

CREATE TABLE Professeurs (
    id_prof     INTEGER PRIMARY KEY,
    nom         TEXT NOT NULL,
    departement TEXT
);

-- Regle 4 : relation N:M -> table d'association
CREATE TABLE Inscriptions (
    id_etudiant INTEGER REFERENCES Etudiants(id_etudiant),
    id_cours    INTEGER REFERENCES Cours(id_cours),
    note        REAL,       -- attribut de la relation
    PRIMARY KEY (id_etudiant, id_cours)
);
\end{minted}
\end{codeblock}
\end{exemple}

\section{Exemple complet~: syst\`eme de biblioth\`eque}

\subsection{Analyse des besoins}
Une biblioth\`eque souhaite g\'erer~:
\begin{itemize}
    \item Les \textbf{livres} (ISBN, titre, ann\'ee de publication).
    \item Les \textbf{auteurs} (nom, nationalit\'e). Un livre peut avoir plusieurs auteurs.
    \item Les \textbf{adh\'erents} (num\'ero, nom, adresse).
    \item Les \textbf{emprunts} (date d'emprunt, date de retour pr\'evue, date de retour effective).
\end{itemize}

\subsection{Diagramme E/R}

\begin{center}
\begin{tikzpicture}[
    entity/.style={rectangle, draw, fill=blue!10, minimum width=2.2cm, minimum height=0.9cm, font=\bfseries\small},
    relationship/.style={diamond, draw, fill=orange!20, minimum width=1.5cm, minimum height=0.8cm, aspect=2, font=\scriptsize\bfseries},
]
    \node[entity] (livre) at (0,0) {Livre};
    \node[entity] (auteur) at (0,-3) {Auteur};
    \node[entity] (adherent) at (7,0) {Adh\'erent};

    \node[relationship] (ecrit) at (0,-1.5) {\'Ecrit};
    \node[relationship] (emprunt) at (3.5,0) {Emprunt};

    \draw (livre) -- node[right] {$N$} (ecrit);
    \draw (ecrit) -- node[right] {$M$} (auteur);
    \draw (livre) -- node[above] {$N$} (emprunt);
    \draw (emprunt) -- node[above] {$M$} (adherent);
\end{tikzpicture}
\end{center}

\subsection{Traduction en SQL}

\begin{codeblock}[title=\textbf{Sch\'ema relationnel de la biblioth\`eque}]
\begin{minted}{sql}
CREATE TABLE Auteurs (
    id_auteur    INTEGER PRIMARY KEY,
    nom          TEXT NOT NULL,
    nationalite  TEXT
);

CREATE TABLE Livres (
    isbn         TEXT PRIMARY KEY,
    titre        TEXT NOT NULL,
    annee_pub    INTEGER
);

CREATE TABLE Ecrit_par (
    isbn         TEXT REFERENCES Livres(isbn),
    id_auteur    INTEGER REFERENCES Auteurs(id_auteur),
    PRIMARY KEY (isbn, id_auteur)
);

CREATE TABLE Adherents (
    id_adherent  INTEGER PRIMARY KEY,
    nom          TEXT NOT NULL,
    adresse      TEXT
);

CREATE TABLE Emprunts (
    id_emprunt   INTEGER PRIMARY KEY,
    isbn         TEXT REFERENCES Livres(isbn),
    id_adherent  INTEGER REFERENCES Adherents(id_adherent),
    date_emprunt DATE NOT NULL,
    date_retour_prevue DATE NOT NULL,
    date_retour_effective DATE
);
\end{minted}
\end{codeblock}

\section{Bonnes pratiques de conception}

\begin{bonnepratique}[title=\textbf{R\`egles d'or de la conception E/R}]
\begin{enumerate}
    \item Chaque type d'entit\'e doit avoir une cl\'e bien d\'efinie.
    \item \'Eviter les attributs multi-valu\'es~: les transformer en entit\'es s\'epar\'ees.
    \item Privil\'egier les relations binaires aux relations ternaires
          (plus faciles \`a comprendre et \`a traduire).
    \item Ne pas confondre entit\'e et attribut~: si un concept a ses propres
          attributs ou participe \`a des relations, c'est une entit\'e.
    \item Documenter chaque choix de conception (cardinalit\'es, contraintes).
\end{enumerate}
\end{bonnepratique}

\begin{antipattern}[title=\textbf{Anti-pattern~: la table ``fourre-tout''}]
Regrouper toutes les informations dans une seule table g\'eante
(\'etudiants, cours, notes, professeurs dans une m\^eme table) entra\^ine~:
\begin{itemize}
    \item Redondance massive des donn\'ees.
    \item Anomalies d'insertion, de mise \`a jour et de suppression.
    \item Difficult\'e de maintenance et d'\'evolution.
\end{itemize}
La normalisation (chapitre~\ref{ch:normalisation}) formalisera ces probl\`emes.
\end{antipattern}

\section{Int\'egration Python~: g\'en\'eration de sch\'ema}

\begin{codeblock}[title=\textbf{G\'en\'eration automatique du sch\'ema avec Python}]
\begin{minted}{python}
import sqlite3

schema = {
    'Auteurs': {
        'id_auteur': 'INTEGER PRIMARY KEY',
        'nom': 'TEXT NOT NULL',
        'nationalite': 'TEXT'
    },
    'Livres': {
        'isbn': 'TEXT PRIMARY KEY',
        'titre': 'TEXT NOT NULL',
        'annee_pub': 'INTEGER'
    }
}

conn = sqlite3.connect('bibliotheque.db')
cur = conn.cursor()

for table, colonnes in schema.items():
    cols = ', '.join(f'{col} {typ}' for col, typ in colonnes.items())
    sql = f'CREATE TABLE IF NOT EXISTS {table} ({cols})'
    print(f'Execution : {sql}')
    cur.execute(sql)

conn.commit()
conn.close()
\end{minted}
\end{codeblock}

\section*{Exercices}

\begin{exercice}
Concevez un diagramme E/R pour un syst\`eme de r\'eservation d'h\^otel
comportant les entit\'es \texttt{Client}, \texttt{Chambre}, \texttt{Hotel}
et la relation \texttt{Reservation}. Pr\'ecisez les cardinalit\'es et
les attributs de chaque entit\'e et relation.
\end{exercice}

\begin{exercice}
Traduisez le diagramme E/R de l'exercice pr\'ec\'edent en sch\'ema relationnel SQL.
Identifiez les cl\'es primaires et \'etrang\`eres.
\end{exercice}

\begin{exercice}
Un h\^opital g\`ere des \texttt{Patients}, des \texttt{M\'edecins}, des
\texttt{Services} et des \texttt{Consultations}. Un patient peut consulter
plusieurs m\'edecins, un m\'edecin appartient \`a un seul service.
\begin{enumerate}
    \item Dessinez le diagramme E/R.
    \item Identifiez une entit\'e faible potentielle.
    \item Traduisez en SQL.
\end{enumerate}
\end{exercice}

\begin{exercice}
Quelle strat\'egie de traduction de sp\'ecialisation choisiriez-vous pour
mod\'eliser un syst\`eme o\`u \texttt{V\'ehicule} est sp\'ecialis\'e en
\texttt{Voiture} et \texttt{Camion}, avec sp\'ecialisation disjointe et totale ?
Justifiez votre choix et donnez le code SQL correspondant.
\end{exercice}

\begin{exercice}
Identifiez et corrigez les erreurs de conception dans le sch\'ema suivant~:
\begin{minted}{sql}
CREATE TABLE Commandes (
    id_commande INTEGER PRIMARY KEY,
    nom_client TEXT,
    adresse_client TEXT,
    telephone_client TEXT,
    produit TEXT,
    prix REAL,
    quantite INTEGER,
    nom_fournisseur TEXT,
    adresse_fournisseur TEXT
);
\end{minted}
Proposez un diagramme E/R correct et le sch\'ema SQL correspondant.
\end{exercice}
